% Options for packages loaded elsewhere
\PassOptionsToPackage{unicode}{hyperref}
\PassOptionsToPackage{hyphens}{url}
\documentclass[
]{article}
\usepackage{xcolor}
\usepackage[margin=1in]{geometry}
\usepackage{amsmath,amssymb}
\setcounter{secnumdepth}{-\maxdimen} % remove section numbering
\usepackage{iftex}
\ifPDFTeX
  \usepackage[T1]{fontenc}
  \usepackage[utf8]{inputenc}
  \usepackage{textcomp} % provide euro and other symbols
\else % if luatex or xetex
  \usepackage{unicode-math} % this also loads fontspec
  \defaultfontfeatures{Scale=MatchLowercase}
  \defaultfontfeatures[\rmfamily]{Ligatures=TeX,Scale=1}
\fi
\usepackage{lmodern}
\ifPDFTeX\else
  % xetex/luatex font selection
\fi
% Use upquote if available, for straight quotes in verbatim environments
\IfFileExists{upquote.sty}{\usepackage{upquote}}{}
\IfFileExists{microtype.sty}{% use microtype if available
  \usepackage[]{microtype}
  \UseMicrotypeSet[protrusion]{basicmath} % disable protrusion for tt fonts
}{}
\makeatletter
\@ifundefined{KOMAClassName}{% if non-KOMA class
  \IfFileExists{parskip.sty}{%
    \usepackage{parskip}
  }{% else
    \setlength{\parindent}{0pt}
    \setlength{\parskip}{6pt plus 2pt minus 1pt}}
}{% if KOMA class
  \KOMAoptions{parskip=half}}
\makeatother
\usepackage{color}
\usepackage{fancyvrb}
\newcommand{\VerbBar}{|}
\newcommand{\VERB}{\Verb[commandchars=\\\{\}]}
\DefineVerbatimEnvironment{Highlighting}{Verbatim}{commandchars=\\\{\}}
% Add ',fontsize=\small' for more characters per line
\usepackage{framed}
\definecolor{shadecolor}{RGB}{248,248,248}
\newenvironment{Shaded}{\begin{snugshade}}{\end{snugshade}}
\newcommand{\AlertTok}[1]{\textcolor[rgb]{0.94,0.16,0.16}{#1}}
\newcommand{\AnnotationTok}[1]{\textcolor[rgb]{0.56,0.35,0.01}{\textbf{\textit{#1}}}}
\newcommand{\AttributeTok}[1]{\textcolor[rgb]{0.13,0.29,0.53}{#1}}
\newcommand{\BaseNTok}[1]{\textcolor[rgb]{0.00,0.00,0.81}{#1}}
\newcommand{\BuiltInTok}[1]{#1}
\newcommand{\CharTok}[1]{\textcolor[rgb]{0.31,0.60,0.02}{#1}}
\newcommand{\CommentTok}[1]{\textcolor[rgb]{0.56,0.35,0.01}{\textit{#1}}}
\newcommand{\CommentVarTok}[1]{\textcolor[rgb]{0.56,0.35,0.01}{\textbf{\textit{#1}}}}
\newcommand{\ConstantTok}[1]{\textcolor[rgb]{0.56,0.35,0.01}{#1}}
\newcommand{\ControlFlowTok}[1]{\textcolor[rgb]{0.13,0.29,0.53}{\textbf{#1}}}
\newcommand{\DataTypeTok}[1]{\textcolor[rgb]{0.13,0.29,0.53}{#1}}
\newcommand{\DecValTok}[1]{\textcolor[rgb]{0.00,0.00,0.81}{#1}}
\newcommand{\DocumentationTok}[1]{\textcolor[rgb]{0.56,0.35,0.01}{\textbf{\textit{#1}}}}
\newcommand{\ErrorTok}[1]{\textcolor[rgb]{0.64,0.00,0.00}{\textbf{#1}}}
\newcommand{\ExtensionTok}[1]{#1}
\newcommand{\FloatTok}[1]{\textcolor[rgb]{0.00,0.00,0.81}{#1}}
\newcommand{\FunctionTok}[1]{\textcolor[rgb]{0.13,0.29,0.53}{\textbf{#1}}}
\newcommand{\ImportTok}[1]{#1}
\newcommand{\InformationTok}[1]{\textcolor[rgb]{0.56,0.35,0.01}{\textbf{\textit{#1}}}}
\newcommand{\KeywordTok}[1]{\textcolor[rgb]{0.13,0.29,0.53}{\textbf{#1}}}
\newcommand{\NormalTok}[1]{#1}
\newcommand{\OperatorTok}[1]{\textcolor[rgb]{0.81,0.36,0.00}{\textbf{#1}}}
\newcommand{\OtherTok}[1]{\textcolor[rgb]{0.56,0.35,0.01}{#1}}
\newcommand{\PreprocessorTok}[1]{\textcolor[rgb]{0.56,0.35,0.01}{\textit{#1}}}
\newcommand{\RegionMarkerTok}[1]{#1}
\newcommand{\SpecialCharTok}[1]{\textcolor[rgb]{0.81,0.36,0.00}{\textbf{#1}}}
\newcommand{\SpecialStringTok}[1]{\textcolor[rgb]{0.31,0.60,0.02}{#1}}
\newcommand{\StringTok}[1]{\textcolor[rgb]{0.31,0.60,0.02}{#1}}
\newcommand{\VariableTok}[1]{\textcolor[rgb]{0.00,0.00,0.00}{#1}}
\newcommand{\VerbatimStringTok}[1]{\textcolor[rgb]{0.31,0.60,0.02}{#1}}
\newcommand{\WarningTok}[1]{\textcolor[rgb]{0.56,0.35,0.01}{\textbf{\textit{#1}}}}
\usepackage{graphicx}
\makeatletter
\newsavebox\pandoc@box
\newcommand*\pandocbounded[1]{% scales image to fit in text height/width
  \sbox\pandoc@box{#1}%
  \Gscale@div\@tempa{\textheight}{\dimexpr\ht\pandoc@box+\dp\pandoc@box\relax}%
  \Gscale@div\@tempb{\linewidth}{\wd\pandoc@box}%
  \ifdim\@tempb\p@<\@tempa\p@\let\@tempa\@tempb\fi% select the smaller of both
  \ifdim\@tempa\p@<\p@\scalebox{\@tempa}{\usebox\pandoc@box}%
  \else\usebox{\pandoc@box}%
  \fi%
}
% Set default figure placement to htbp
\def\fps@figure{htbp}
\makeatother
\setlength{\emergencystretch}{3em} % prevent overfull lines
\providecommand{\tightlist}{%
  \setlength{\itemsep}{0pt}\setlength{\parskip}{0pt}}
\usepackage{bookmark}
\IfFileExists{xurl.sty}{\usepackage{xurl}}{} % add URL line breaks if available
\urlstyle{same}
\hypersetup{
  pdftitle={r-import lesson},
  pdfauthor={Vanessa},
  hidelinks,
  pdfcreator={LaTeX via pandoc}}

\title{r-import lesson}
\author{Vanessa}
\date{2026-09-21}

\begin{document}
\maketitle

\subsection{R Markdown}\label{r-markdown}

This is an R Markdown document. Markdown is a simple formatting syntax
for authoring HTML, PDF, and MS Word documents. For more details on
using R Markdown see \url{http://rmarkdown.rstudio.com}.

When you click the \textbf{Knit} button a document will be generated
that includes both content as well as the output of any embedded R code
chunks within the document. You can embed an R code chunk like this:

\begin{Shaded}
\begin{Highlighting}[]
\FunctionTok{summary}\NormalTok{(cars)}
\end{Highlighting}
\end{Shaded}

\begin{verbatim}
##      speed           dist       
##  Min.   : 4.0   Min.   :  2.00  
##  1st Qu.:12.0   1st Qu.: 26.00  
##  Median :15.0   Median : 36.00  
##  Mean   :15.4   Mean   : 42.98  
##  3rd Qu.:19.0   3rd Qu.: 56.00  
##  Max.   :25.0   Max.   :120.00
\end{verbatim}

\subsection{Including Plots}\label{including-plots}

You can also embed plots, for example:

\pandocbounded{\includegraphics[keepaspectratio]{r-import_files/r-import_files/figure-latex/pressure-1.pdf}}

dir.create(``data'') dir.create(``data/linelists'')
install.packages(``rio'') \# universal import/export tool
install.packages(``here'') \# builds reliable file paths
install.packages(``pacman'') \# lets you load multiple packages in one
line (used later) library(rio) library(here) here() linelist \textless-
import(here(``data'', ``linelists'', ``linelist\_cleaned.xlsx''))
head(linelist) gss \textless- import(here(``data'', ``GSSsubset.csv''))
head(gss) my\_data \textless- import(here(``data'', ``linelists'',
``linelist\_cleaned.xlsx''), which = ``Sheet1'') my\_data \textless-
import(here(``data'', ``linelists'', ``linelist\_cleaned.xlsx''), which
= ``linelist'') readxl::excel\_sheets(here(``data'', ``linelists'',
``lmy\_data \textless- import(here(''data'', ``linelists'',
``linelist\_cleaned.xlsx''), which = ``Sheet
1'')inelist\_cleaned.xlsx'')) head(my\_data) linelist \textless-
import(here(``data'', ``linelists'', ``linelist\_cleaned.xlsx''), na =
c(``99'', ``Missing'', ``\,``)) linelist\_raw \textless-
import(here(''data'', ``linelists'', ``linelist\_cleaned.xlsx''), skip =
1) export(linelist, here(``data'', ``clean\_linelist.xlsx'')) \# save as
Excel export(linelist, here(``data'', ``clean\_linelist.csv'')) \# save
as CSV export(linelist, here(``data'', ``clean\_linelist.rds'')) \# save
as R's own format linelist2 \textless- import(here(``data'',
``clean\_linelist.rds'')) head(linelist2) linelist2 \textless-
import(here(``data'', ``clean\_linelist.rds''), trust = FALSE) \# create
a fake messy file to practice on messy \textless- data.frame( age =
c(``Age in years'', ``34'', ``51'', ``22''), sex = c(``Patient sex'',
``Male'', ``Female'', ``Male'') ) export(messy, here(``data'',
``messy\_example.csv'')) view(messy) bad\_import \textless-
import(here(``data'', ``messy\_example.csv'')) bad\_import
library(dplyr) \# 1. Import once just to grab the real column names
real\_names \textless- import(here(``data'', ``messy\_example.csv''))
\%\textgreater\% names() real\_names

\section{2. Import again, skipping the messy second row, then reapply
names}\label{import-again-skipping-the-messy-second-row-then-reapply-names}

clean\_import \textless- import(here(``data'', ``messy\_example.csv''),
skip = 1, col.names = real\_names) clean\_import library(tibble)

manual\_entry \textless- tribble( \textasciitilde patient\_id,
\textasciitilde age, \textasciitilde outcome, 1, 34, ``Recovered'', 2,
51, ``Died'', 3, 22, ``Recovered'' ) manual\_entry patient\_id
\textless- c(1, 2, 3) age \textless- c(34, 51, 22) outcome \textless-
c(``Recovered'', ``Died'', ``Recovered'')

manual\_entry2 \textless- data.frame(patient\_id, age, outcome)
manual\_entry2 install.packages(``clipr'') library(clipr)

pasted\_data \textless- read\_clip\_tbl( HBA1C Lipids CRP CBC KFTs Uric
acids Urinalysis Spirometry respiratory genetics DKL 2190 2690 1990 990
3390 790 890\\
pathcare 2640 ~3,700 2,490 1,480 3,070 ~840\\
metropolis 2425 2625 2475 1300 2975 825 975\\
Crystal lab 1000\\
) pasted\_data \textless- read\_clip\_tbl() pasted\_data
view(pasted\_data)

\end{document}
